\section{Results}
\label{sec:results}

Our experiments validate the core IPCR mechanism while revealing important limitations. We present results for each research question, with interpretation of what the findings mean for our hypothesis.

\subsection{Main Results: IPCR Achieves 95\% of Oracle Performance}

Table~\ref{tab:main} presents end-to-end routing performance on held-out FLAN task families.

\begin{table}[h]
\centering
\begin{tabular}{lcc}
\toprule
\textbf{Routing Strategy} & \textbf{Performance} & \textbf{Relative to Oracle} \\
\midrule
Oracle & 91.39\% & 100.00\% \\
\textbf{IPCR (ours)} & \textbf{86.82\%} & \textbf{95.00\%} \\
Uniform & 36.79\% & 40.26\% \\
Random & 14.97\% & 16.38\% \\
\bottomrule
\end{tabular}
\caption{End-to-end routing performance on FLAN task families.}
\label{tab:main}
\end{table}

\textbf{Key finding:} IPCR achieves 95.00\% of oracle performance, exceeding our 90\% threshold. The improvement over Uniform is highly significant ($t = 68.99$, $p = 7.05 \times 10^{-224}$), confirming that routing provides substantial benefit beyond naive adapter aggregation.

This result validates our central claim: zero-shot adapter routing is achievable without validation data when instruction embeddings contain sufficient task signal.

\subsection{Task Family Separability (H-E0)}

The prerequisite for effective routing is that instruction embeddings cluster by task type.

\begin{table}[h]
\centering
\begin{tabular}{lccc}
\toprule
\textbf{Metric} & \textbf{Value} & \textbf{Threshold} & \textbf{Status} \\
\midrule
Macro-F1 & 0.995 & $\geq$0.75 & \cmark~Pass \\
Accuracy & 0.998 & --- & --- \\
Baseline F1 & 0.115 & --- & --- \\
\bottomrule
\end{tabular}
\caption{Task family separability results.}
\label{tab:separability}
\end{table}

\textbf{Per-family F1 scores:} All 9 tested families exceed F1 = 0.98, ranging from 0.982 (strategyqa) to 1.000 (sensemaking).

\textbf{Interpretation:} Near-perfect linear separability (F1 = 0.995) demonstrates that MiniLM embeddings encode strong task-type signal. This aligns with our hypothesis that instruction semantics and adapter specializations share geometric structure. The 8.6$\times$ improvement over random baseline (0.995 vs 0.115) confirms the signal is genuine, not artifacts.

Dimensionality reduction (t-SNE) of the embedding space reveals clear clustering by task family, with minimal overlap between clusters (see supplementary materials for visualization). The per-family F1 scores (0.982--1.000) confirm that clusters are not only visually distinct but linearly separable.

\subsection{Adapter Selection Accuracy (H-E1)}

Given separable embeddings, can the linear probe select the correct adapter?

\begin{table}[h]
\centering
\begin{tabular}{lccc}
\toprule
\textbf{Metric} & \textbf{Value} & \textbf{Threshold} & \textbf{Status} \\
\midrule
Top-1 Accuracy & 72.67\% & $\geq$70\% & \cmark~Pass \\
Top-3 Accuracy & 95.78\% & $\geq$85\% & \cmark~Pass \\
vs Random & 14.8$\times$ & $>$1$\times$ & \cmark~Pass \\
\bottomrule
\end{tabular}
\caption{Adapter selection accuracy.}
\label{tab:selection}
\end{table}

\textbf{Interpretation:} The high top-3 accuracy (95.78\%) indicates that even when top-1 prediction is wrong, the correct adapter is nearly always in the top 3 candidates. This supports a routing-with-fallback strategy for deployment: select top-1 by default, but top-3 coverage provides safety margin.

The 14.8$\times$ improvement over random selection confirms the linear probe extracts meaningful routing signal from frozen embeddings.

\subsection{Robustness Analysis (H-M2): A Critical Limitation}

We tested whether routing is robust to instruction variations---a requirement for real-world deployment.

\subsubsection{Paraphrase Robustness}

\begin{table}[h]
\centering
\begin{tabular}{lcc}
\toprule
\textbf{Variant} & \textbf{Cosine Mean} & \textbf{Routing Consistency} \\
\midrule
WordNet synonyms & 0.720 & 72.5\% \\
Embedding-filtered & 0.843 & 79.6\% \\
\textbf{Combined} & \textbf{0.782} & \textbf{76.1\%} \\
\bottomrule
\end{tabular}
\caption{Paraphrase robustness results.}
\label{tab:paraphrase}
\end{table}

\textbf{Threshold:} Cosine $\geq$ 0.90. \textbf{Result:} \xmark~Fail (0.782)

\subsubsection{Keyword Masking}

\begin{table}[h]
\centering
\begin{tabular}{lccc}
\toprule
\textbf{Masking Rate} & \textbf{Original Acc} & \textbf{Masked Acc} & \textbf{Drop} \\
\midrule
20\% keywords & 74.9\% & 48.9\% & 26.0\% \\
50\% keywords & 74.9\% & 30.4\% & 44.4\% \\
\bottomrule
\end{tabular}
\caption{Keyword masking results.}
\label{tab:masking}
\end{table}

\textbf{Threshold:} Drop $<$ 10\%. \textbf{Result:} \xmark~Fail (44.4\%)

\textbf{Interpretation:} The robustness failure reveals a fundamental characteristic of IPCR: routing depends on lexical keywords, not deep semantic invariants. WordNet synonym substitutions (e.g., ``calculate'' $\rightarrow$ ``compute'', ``capital'' $\rightarrow$ ``Washington'') cause embedding drift beyond routing tolerance.

This is not a minor limitation---it determines IPCR's deployment scope. The method is appropriate for controlled instruction formats (APIs, templates, chatbot interfaces with canonical phrasing) but fragile for open-ended queries with arbitrary paraphrasing.

\textbf{Root cause analysis:} MiniLM's mean-pooling architecture captures token presence rather than compositional meaning. Task discrimination emerges from keyword clusters in FLAN templates, not from understanding instruction semantics. This ``lexical anchoring'' hypothesis explains both our success (FLAN instructions contain task-indicative keywords) and our limitation (paraphrased instructions lose the anchor tokens).

\subsection{Per-Class Analysis}

Robustness varies substantially across task families:

\begin{table}[h]
\centering
\begin{tabular}{lll}
\toprule
\textbf{Task Family} & \textbf{Routing Consistency} & \textbf{Interpretation} \\
\midrule
cot\_qasc & 97.2\% & Robust --- distinctive vocabulary \\
stream\_aqua & 89.0\% & Robust --- math keywords \\
cot\_esnli & 49.2\% & Fragile --- open-domain NLI \\
cot\_creak & 62.9\% & Fragile --- claim verification \\
\bottomrule
\end{tabular}
\caption{Per-class robustness analysis.}
\label{tab:perclass}
\end{table}

Tasks with distinctive vocabulary (math: ``calculate'', ``equation'') maintain routing stability. Open-domain tasks with less keyword anchoring are most vulnerable.
